% ===== PRÉAMBULE CANONIQUE — à coller VERBATIM en tête de CHAQUE .tex =====
\documentclass[11pt,a4paper]{article}
\usepackage[utf8]{inputenc}\usepackage[T1]{fontenc}\usepackage{lmodern}
\usepackage[french]{babel}
\usepackage{amsmath,amssymb,mathtools}\usepackage{icomma}
\usepackage[version=4]{mhchem}          % \ce{...} : équations & formules
\usepackage{siunitx}\sisetup{locale=FR} % \qty{}{}, \si{}, \ang{}
\usepackage{chemfig}                    % \chemfig{...} : structures développées
\usepackage{xcolor}\usepackage{colortbl}
\usepackage{tikz}\usepackage{pgfplots}\pgfplotsset{compat=1.18}
\usepackage[most]{tcolorbox}\usepackage{enumitem}\usepackage{array,booktabs}
\usepackage[a4paper,margin=2.2cm]{geometry}
\usetikzlibrary{arrows.meta,calc,angles,quotes,positioning,patterns,backgrounds,shapes.geometric}
\definecolor{primary}{RGB}{222,49,99}\definecolor{primarydark}{RGB}{150,22,55}
\definecolor{defcol}{RGB}{0,114,178}\definecolor{methcol}{RGB}{52,92,148}
\definecolor{excol}{RGB}{0,135,90}\definecolor{attncol}{RGB}{200,40,55}
\tcbset{boxbase/.style={enhanced,breakable,boxrule=0.9pt,arc=2.5pt,
  left=3mm,right=3mm,top=2.3mm,bottom=2.3mm,fonttitle=\bfseries,coltitle=white}}
% --- boîtes TUTORIEL (noms EXACTS, ne pas renommer) ---
\newtcolorbox{cle}[1][]{boxbase,colback=defcol!6,colframe=defcol,
  title={\textsc{À retenir}\ifx\relax#1\relax\relax\else\textmd{ : \itshape #1}\fi}}
\newtcolorbox{methode}[1][]{boxbase,colback=methcol!7,colframe=methcol,sharp corners,
  title={\textsc{Méthode}\ifx\relax#1\relax\relax\else\textmd{ --- \itshape #1}\fi}}
\newtcolorbox{exemplebox}[1][]{boxbase,colback=excol!5,colframe=excol,
  title={\textsc{Exemple}\ifx\relax#1\relax\relax\else\textmd{ : \itshape #1}\fi}}
\newtcolorbox{piege}[1][]{boxbase,colback=attncol!6,colframe=attncol,
  title={\textsc{Piège}\ifx\relax#1\relax\relax\else\textmd{ : \itshape #1}\fi}}
\newtcolorbox{remarque}[1][]{boxbase,colback=black!4,colframe=black!45,
  title={\textsc{Remarque}\ifx\relax#1\relax\relax\else\textmd{ : \itshape #1}\fi}}
% --- boîtes EXERCICE / CORRIGÉ (noms EXACTS, auto-comptées) ---
\newtcolorbox[auto counter]{exo}[1][]{boxbase,colback=primary!4,colframe=primary,
  title={\textsc{Exercice}~\thetcbcounter\ifx\relax#1\relax\relax\else\textmd{ --- \itshape #1}\fi}}
\newtcolorbox[auto counter]{corrige}[1][]{boxbase,colback=excol!4,colframe=excol,
  title={\textsc{Corrigé}~\thetcbcounter\ifx\relax#1\relax\relax\else\textmd{ --- \itshape #1}\fi}}
\newcommand{\R}{\mathbb{R}}\newcommand{\N}{\mathbb{N}}\newcommand{\Z}{\mathbb{Z}}\newcommand{\ds}{\displaystyle}
% (un \usepackage{fancyhdr}+en-tête est possible ; il est ignoré côté web)
% =========================================================================
\begin{document}
\begin{center}\color{primarydark}\large\bfseries Thème 2 — Corrigé \quad$\bullet$\quad Niveau Difficile\end{center}

\begin{corrige}[dichromate et acide nitreux]
Réduction : \ce{Cr2O7^{2-} + 14 H+ + 6 e^- <=> 2 Cr^{3+} + 7 H2O} ($2\times 3=6$ électrons).
Oxydation : \ce{HNO2 + H2O <=> NO3^- + 3 H+ + 2 e^-} ($n=5-3=2$). On multiplie l'oxydation par 3
(PPCM $=6$), on additionne et l'on simplifie \ce{H+} ($14-9$) et \ce{H2O} ($7-3$) :
\[
\ce{Cr2O7^{2-} + 3 HNO2 + 5 H+ -> 2 Cr^{3+} + 3 NO3^- + 4 H2O}
\]
Contrôle : $\ce{O}:7+6=13=4+9$ \checkmark ; charge : $-2+5=+3=2(+3)-3$ \checkmark. Il s'échange
\textbf{6 électrons}.
\end{corrige}

\begin{corrige}[la réaction aluminothermique]
\begin{enumerate}[label=\alph*)]
  \item \ce{2 Al + Fe2O3 -> Al2O3 + 2 Fe}.
  \item L'aluminium passe de $0$ à $+\mathrm{III}$ : il cède \textbf{3 électrons} par atome (au total
        $2\times 3=6$ électrons, captés par les 2 \ce{Fe^{3+}}).
  \item Coefficients $2,1,1,2$ : leur somme vaut $\mathbf{6}$. L'affirmation \textbf{(C)} (« vaut 4 »)
        est donc \textbf{incorrecte}. (A, B et D sont exactes.)
\end{enumerate}
\end{corrige}

\begin{corrige}[oxydation en milieu basique]
Demi-équations en milieu basique :
\[
\ce{MnO4^- + 2 H2O + 3 e^- <=> MnO2 + 4 OH^-}\qquad
\ce{SO3^{2-} + 2 OH^- <=> SO4^{2-} + H2O + 2 e^-}
\]
PPCM des électrons $=6$ : on multiplie la première par 2, la seconde par 3, on additionne et l'on
simplifie \ce{H2O} et \ce{OH^-} communs :
\[
\ce{2 MnO4^- + 3 SO3^{2-} + H2O -> 2 MnO2 + 3 SO4^{2-} + 2 OH^-}
\]
Contrôle : $\ce{O}:8+9+1=18=4+12+2$ \checkmark ; charge : $-2-6=-8=-6-2$ \checkmark.
\end{corrige}

\begin{corrige}[permanganate et iodure]
Réduction : \ce{MnO4^- + 8 H+ + 5 e^- <=> Mn^{2+} + 4 H2O}. Oxydation : \ce{2 I^- <=> I2 + 2 e^-}.
PPCM $=10$ : on multiplie la réduction par 2 et l'oxydation par 5.
\begin{enumerate}[label=\alph*)]
  \item \[\ce{2 MnO4^- + 10 I^- + 16 H+ -> 2 Mn^{2+} + 5 I2 + 8 H2O}\]
        (charge : $-2-10+16=+4=2\times(+2)$ \checkmark).
  \item Il s'échange \textbf{10 électrons} au total.
\end{enumerate}
\end{corrige}

\begin{corrige}[le cuivre dans l'acide nitrique : dilué ou concentré]
Le cuivre s'oxyde ($\ce{Cu <=> Cu^{2+} + 2 e^-}$) dans les deux cas ; seule change la réduction de
l'azote.
\begin{enumerate}[label=\alph*)]
  \item \ce{NO3^-}$\to$\ce{NO} : $n=5-2=3$. PPCM$(2,3)=6$ :
        \[\ce{3 Cu + 2 NO3^- + 8 H+ -> 3 Cu^{2+} + 2 NO + 4 H2O}\]
  \item \ce{NO3^-}$\to$\ce{NO2} : $n=5-4=1$. PPCM$(2,1)=2$ :
        \[\ce{Cu + 2 NO3^- + 4 H+ -> Cu^{2+} + 2 NO2 + 2 H2O}\]
\end{enumerate}
\end{corrige}

\end{document}
